\chapter{살아 있는 세포로 만든 컴퓨터}
% full version: chapters/21_living.tex

\pencilfig{21}{\pencilart{21}}{칩 위의 살아 있는 뉴런. 그 아래에 전극 수천 개가 있고, 전극마다 듣기도 하고 말하기도 한다.}

실리콘이 아니라 살아 있는 뉴런으로 기계를 만들 수 있을까? 만들 수 있고, 이미 만들고 있다. 전극이 수천 개 달린 칩 위에서 바로 뉴런을 기른다. 평평한 층으로 기르기도 하고, 작은 뇌를 닮은 덩어리로 기르기도 한다. 전극은 저마다 듣기도 하고 신호를 넣기도 한다. 회로가 열려 있는 테스트 벤치다. 모든 것을 측정할 수 있고 모든 것에 개입할 수 있다. 다만 방송국은 없다. 제어 레지스터 없는 데이터 버스다.

\section*{네트워크가 스스로 하는 일}

그대로 두면 배양 전체가 이따금 한꺼번에 불붙는다. 거의 모든 뉴런이 함께 딸깍한 뒤 조용해지고, 그 간격은 1초에서 몇 분까지다. 이미 익숙한 발진기다. 네트워크는 스스로를 흥분시키다가 지치고, 연결이 물질을 다시 채울 때까지 쉰다.

\pencilfig{129}{%
% spike raster of 8 neurons in a culture left to itself: sparse single clicks and shared bursts
\foreach \b in {0.9,3.1,4.0,6.6}{\fill[pcAmber, opacity=0.18] (\b-0.1,-0.15) rectangle (\b+0.32,2.95);}
\foreach \y/\ss in {0/{0.96,1.0,1.13,2.43,3.0,3.12,3.24,4.02,4.09,4.15,6.66,6.69,6.81},
  1/{0.46,0.88,1.0,1.05,3.09,3.2,3.94,4.04,4.37,6.65,6.71},
  2/{0.73,0.84,0.94,1.41,3.08,3.12,3.21,4.05,4.06,4.17,6.59,6.63},
  3/{0.89,0.93,1.06,3.1,3.19,3.28,3.89,3.94,4.13,4.2,4.31,6.63,6.65,6.78},
  4/{0.87,0.96,3.09,3.15,3.23,3.73,3.96,4.06,4.19,5.98,6.66,6.68,6.75},
  5/{0.44,0.92,0.97,3.05,3.22,3.27,4.01,4.07,4.17,5.76,6.57,6.73,6.75},
  6/{0.92,1.0,1.73,2.85,3.18,3.19,4.0,4.07,6.53,6.66,6.73},
  7/{0.87,0.92,3.13,3.13,3.27,3.99,4.06,4.17,6.44,6.61,6.72,7.13}}{
  \foreach \s in \ss {\draw[pcBlue, line width=0.7pt] (\s,0.35*\y) -- (\s,0.35*\y+0.25);}}
\draw[arr] (0,-0.3) -- (7.5,-0.3); \node[lbl, anchor=north] at (3.75,-0.32) {시간};
\node[lbl, anchor=east, align=right] at (-0.05,1.35) {뉴런\\8개};
\node[lbl, text=pcAmber!70!black, anchor=south] at (3.55,2.95) {공동 분출};
}{그대로 둔 칩 위의 배양. 따로 떨어진 딸깍은 드물고, 이따금 네트워크 거의 전체가 한꺼번에 불붙었다가 연결이 물질을 다시 채울 때까지 조용해진다.}

\section*{도파민 없는 교사}

도파민이 없다면 이런 네트워크를 어떻게 가르칠까? 처음 쓰인 방법은 전기적인 것이었다. 예를 들면, 네트워크가 원하는 응답을 낼 때까지 자극하다가 그 순간 멈춘다. 그러면 그 응답이 굳어진다. 다른 실험에서는 배양이 간단한 가상의 “몸”을 조종하며 그것을 목표로 이끄는 법을 배웠다. 이 모든 실험에서 교사는 엔지니어가 고른 전류 패턴이다.

\section*{빛의 섬광으로 푸는 도파민}

한번은 교사를 화학적으로 만들어 보려는 시도가 있었다. 살아 있는 뉴런 덩어리를 쓰는 한 플랫폼에서, 배양액에 빛에 민감한 “우리”에 갇혀 작용하지 못하는 도파민을 넣었다. 네트워크가 자극에 전보다 강하게 응답하면 “상”을 주었다. 자외선 섬광이 우리를 부수면 도파민이 풀려나는 것이다. 240번 반복하는 동안 응답은 전체적으로 커졌다. 그러나 저자들 스스로 이런 실험 하나로는 그 원인이 도파민인지 말할 수 없다고 솔직하게 쓰고 있다.

\section*{수레 위의 막대}

최근의 한 실험에서는 뉴런 덩어리에게 수레 위의 막대를 세워 균형을 잡도록 가르쳤다. 학습 프로그램의 고전적인 과제다. 어떤 학습 펄스를 넣을지는 프로그램이 골랐다. 프로그램이 고른 신호가 무작위 신호보다 잘 들었지만, 45분 쉬고 나면 배운 것이 사라졌다. 그리고 \nt{GLUT} 연결을 막으면 학습이 전혀 일어나지 않았다. 배운 것은 바로 거기에 산다. 하지만 교사는 여전히 바깥에 있다. 다만 이제는 엔지니어가 아니라 프로그램이다.

\section*{살아 있는 저수조}

다른 접근도 있다. 살아 있는 네트워크를 아예 가르치지 않는 것이다. 네트워크를 “저수조”로 쓴다. 풍부하고 얽히고설킨 시스템이고, 단순한 전자 판독기가 그 응답을 풀어 읽는 법을 배운다. 이런 식으로 살아 있는 덩어리가 목소리를 구별하는 일을 도왔다. 다만 저자들에 따르면 작업하는 동안 덩어리 안의 연결도 바뀌었다.

\section*{비용과 질문}

칩 위의 뉴런 백만 개는 약 1만분의 1와트를 쓴다. 주변의 전자 장치보다 훨씬 적다. 그리고 이 테스트 벤치는 무엇을 성공으로 칠지 스스로 정하지 못한다. 그 역할은 언제나 바깥의 누군가가 맡는다. 이런 시스템이 커질수록 윤리의 문제도 떠오르며, 이 문제는 아직 아무도 풀지 못했다.

\fullversion{21}
